\documentclass[10pt,a4paper]{amsart}
\usepackage[margin=19mm]{geometry}
\usepackage{amsmath,amssymb,mathtools}
\usepackage[scr=rsfs,cal=euler]{mathalfa}
\usepackage{xfrac}
\usepackage[unicode,hidelinks,bookmarks=false]{hyperref}
\newcommand{\ie}{i.e.\ }
\newcommand{\cf}{cf.\ }
\newcommand{\resp}{respectively\ }
\newcommand{\Subsection}{\subsection}
\providecommand{\nobreakdash}{\nobreak\hbox{-}}
\setlength{\emergencystretch}{2em}
\multlinegap0pt
\usepackage{bm}
\usepackage{bbm}
\usepackage{dsfont}
\usepackage{xspace}
\usepackage{cancel}
\usepackage{mathtools}
\usepackage[shortlabels]{enumitem}
\usepackage{amsthm}
\makeatletter
\@ifundefined{thm}{\newtheorem{thm}{Theorem}}{}
\@ifundefined{lem}{\newtheorem{lem}[thm]{Lemma}}{}
\@ifundefined{prop}{\newtheorem{prop}[thm]{Proposition}}{}
\makeatother
\DeclareMathOperator{\Ball}{Ball}
\DeclareMathOperator{\Lip}{Lip}
\DeclareMathOperator{\md}{md}
\DeclarePairedDelimiter{\ip}{\langle}{\rangle}
\newcommand{\bC}{{\mathbb C}}
\newcommand{\bM}{{\mathbb M}}
\newcommand{\bN}{{\mathbb N}}
\newcommand{\cU}{{\mathcal U}}
\title[On Generators for $W^{*}$-bundles]{On Generators for $W^{*}$-bundles}
\author{Ben Hayes}
\date{Source: arXiv:2609.16428v2 (CC BY 4.0)}
\begin{document}
\maketitle

\noindent\emph{Source and licence.} On Generators for $W^{*}$-bundles. Version arXiv:2609.16428v2 (CC BY 4.0), distributed under the Creative Commons Attribution 4.0 International licence, as recorded in the arXiv licence metadata for this exact version and shown on the abstract page https://arxiv.org/abs/2609.16428v2. Full rights notice: licenses/arxiv-2609.16428-CC-BY-4.0.txt.\par\smallskip

\noindent\textit{Editorial scope.} Counted unit: the Thm whose source label is \texttt{thm: invariance of MD} in the article (source0.tex lines 600--607, source label \texttt{thm: invariance of MD}), delivered together with the article's own complete original proof of it (source0.tex lines 610--659, one \texttt{proof} environment of 50 lines). No mathematics is written, repaired or re-derived: statement and proof are the article's own text line for line. Required context retained uncounted: prop: better KDT (lines 554--556); lem: key lem for invariance (lines 561--568); item: pushforward neighborhods (lines 563--567); item: small perturbation neighborhood (lines 563--567); item: small perturbation neighborhood 2 (lines 563--567). Every definition, display and label of those lines is retained unchanged. Omitted from this package and not redistributed: 1--553, 557--560, 568--599, 608--609, 660--2306. Nothing inside a retained excerpt is omitted or altered: every retained body is delivered exactly as the article's lines are written, so no omission receipt of any kind accompanies this package. Citations: the retained excerpts carry no citation command at all, so no bibliography entry of the article is transcribed and no reference is redistributed. Reference adaptation: none was needed and none was made; every reference inside the delivered unit resolves inside it, so no unresolved reference or citation is produced. Printed slips kept literal: no printed word, symbol, hypothesis or number of the counted statement or of its proof was altered. Layout and class: the article's own class is \texttt{amsart} with packages including amsmath, amscd, amssymb, amsthm, enumerate, ulem, mathtools, xcolor, stmaryrd, geometry, fontenc, inputenc, mathrsfs; the wrapper is the lane's pinned \texttt{amsart} editorial wrapper at 10pt on A4, which restates the theorem-like environments the retained text needs and adds the article's own macro definitions that the retained text needs (\texttt{\textbackslash Ball}, \texttt{\textbackslash Lip}, \texttt{\textbackslash bC}, \texttt{\textbackslash bM}, \texttt{\textbackslash bN}, \texttt{\textbackslash cU}, \texttt{\textbackslash ip}, \texttt{\textbackslash md}), with extra wrapper packages (bm, bbm, dsfont, xspace, cancel, mathtools, enumitem). The delivered numbering is the unit's own. Assets: no retained span contains a figure environment, a graphics-inclusion command, a PDF-page inclusion, a table or a TikZ picture; no image file is copied into this package, the package declares no asset dependency and no image file is redistributed with it. Licence: version arXiv:2609.16428v2 of this article is distributed under the Creative Commons Attribution 4.0 International licence, as recorded in the arXiv licence metadata for this exact version and shown on the abstract page https://arxiv.org/abs/2609.16428v2. A full rights notice is delivered at licenses/arxiv-2609.16428-CC-BY-4.0.txt. Characters: every retained excerpt is pure ASCII, so the delivered engine, XeLaTeX with its default Latin Modern text font, reports no missing character.
\begin{prop}\label{prop: better KDT}
Let $(M,X)$ be a tracially complete $C^{*}$-algebra, $I$ a set and $a\in M^{I}$. Let $R\in [0,\infty)^{I}$ be a cutoff function for $a$. For any $b\in C^{*}_{X}(a)$, and any $\varepsilon>0$, there is a $P\in \bC^{*}\ip{(T_{i})_{i\in I}}$ so that $\|P\|_{R}\leq \|b\|$ and $\|P(a)-b\|_{2;X}<\varepsilon$. 
\end{prop}

\begin{lem}\label{lem: key lem for invariance}
Let $(M,X)$ be a tracially complete $C^{*}$-algebra, and let $I$ be a set. Fix $x\in M^{I}$, and a cutoff function $R\in [0,\infty)^{I}$. 
\begin{enumerate}[(i)]
\item Let $Q\in \bC^{*}\ip{T_{i}:i\in I}^{I'}$ for some set $I'$ and let $R'\in [0,+\infty)^{I'}$ satisfy that $\|Q_{i'}\|_{R,\infty}\leq R_{i}'$ for all $i'\in I'$.Given a weak$^{*}$-neighborhood $\mathcal{O}$ of $L_{Q(x);X}$, there is a weak$^{*}$-neighborhood $\mathcal{V}$ of $L_{x;X}$ so that $Q(\Gamma^{(n)}_{R}(\mathcal{V}))\subseteq \Gamma^{(n)}_{R'}(\mathcal{O})$ for all $n\in \bN$. \label{item: pushforward neighborhods}
\item For any weak$^{*}$-neighborhood $\mathcal{O}$ of $L_{x;X}$, there is a $\kappa>0$, a finite $F\subseteq I$, and a weak$^{*}$-neighborhood $\mathcal{V}$ of $L_{x;X}$ so that $N_{\kappa}(\Gamma^{(n)}(\mathcal{V}),\|\cdot\|_{2,F}))\cap \left(\prod_{i\in I}R_{i}\Ball(\bM_{n}(\bC))\right)\subseteq \Gamma^{(n)}_{R}(\mathcal{O})$ for all $n\in \bN$. \label{item: small perturbation neighborhood}
\item For any weak$^{*}$-neighborhood $\mathcal{O}$ of $L_{x;X}$, there is an $\eta>0$ and a finite $F\subseteq I$ with the following property. Whenever $b\in \prod_{i\in I}R_{i}\Ball(M)$ satisfies $\|x-b\|_{2,F;X}<\eta$, there is a weak$^{*}$-neighborhood $\mathcal{V}$ of $L_{b;X}$ with $\Gamma^{(n)}_{R}(\mathcal{V})\subseteq \Gamma^{(n)}_{R}(\mathcal{O})$.  \label{item: small perturbation neighborhood 2}
\end{enumerate}
\end{lem}

\begin{thm}\label{thm: invariance of MD}
Let $(M,X)$ be a tracially complete $C^{*}$-algebra, and $(I_{i})_{i=1}^{3},(I_{i}')_{i=1}^{3}$ be sets. Suppose that $a\in M^{I_{1}},b\in M^{I_{2}},c\in M^{I_{3}}$, and $a'\in M^{I_{1}'},b'\in M^{I_{2}'},c'\in M^{I_{3}'}$ are given, and that 
\[C^{*}_{X}(b)\supseteq C^{*}_{X}(b'),\]
\[C^{*}_{X}(a,b)\subseteq C^{*}_{X}(a',b'),\]
\[C^{*}_{X}(a,b,c)\supseteq C^{*}_{X}(a',b',c').\]
Then
\[\delta_{MD}(a|b:c)\leq\delta_{MD}(a'|b':c').\]
\end{thm}

\begin{proof}
To ease notation somewhat, set $I=\sqcup_{i=1}^{3}I_{i}$ and $I'=\sqcup_{i=1}^{3}I_{i}'$.
 Fix cutoff functions $R\in [0,\infty)^{I}$ and $R'\in [0,\infty)^{I'}$. Let $\varepsilon>0$  and a finite $F\subseteq I_{1}\sqcup I_{2}$ be given. 
 Since $C^{*}_{X}(a,b)\subseteq C^{*}_{X}(a',b')$, we may apply Proposition \ref{prop: better KDT} to choose a $P\in \bC^{*}\ip{T_{i}':i\in I_{1}'\sqcup I_{2}'}^{I_{1}\sqcup I_{2}}$ so that $\|P(a',b')-(a,b)\|_{2,F;X}<\varepsilon$ and that $\|P_{i}\|_{R'}\leq R_{i}$ for all $i\in I_{1}\sqcup I_{2}$. Let $F'\subseteq I_{1}'\sqcup I_{2}'$ be a finite set so that $P_{i}\in \bC^{*}\ip{T_{i}':i\in F'}$ for all $i\in F$.  Set $L'=\|P|_{F}\|_{R',\Lip}$.  Set $\varepsilon'=\frac{\varepsilon}{L'+1}$. Since $C^{*}_{X}(a',b',c')\subseteq C^{*}_{X}(a,b,c)$, we may apply Proposition \ref{prop: better KDT} to find a $Q_{0}\in \bC^{*}\ip{T_{i'}:i'\in I'}^{I}$ so that 
 \[\|Q_{0}(a,b,c)-(a',b',c')\|_{2,F'}<\varepsilon'.\]
Note that
 \[\|P(Q_{0}|_{I_{1}'\sqcup I_{2}'}(a,b,c))-(a,b)\|_{2,F;X}<2\varepsilon.\]
 Thus 
 \[\cU_{1}=\left\{\ell\in \Sigma_{R}:\sum_{i\in F}\|P_{i}(Q_{0}|_{I_{1}'\sqcup I_{2}'}(T'))-T_{i}\|_{L^{2}(\ell)}^{2}<4\varepsilon^{2}\right\}\]
 is a weak$^{*}$-neighborhood of $L_{a,b,c;X}$. 

  Fix a weak$^{*}$-open neighborhood $\mathcal{O}$ of $L_{a',b',c';X}$, a finite $G'\subseteq I_{2}'$, and a $\theta'>0$. By Lemma \ref{lem: key lem for invariance} (\ref{item: small perturbation neighborhood}) we may find an $\eta>0$ and a finite $\widetilde{F}\subseteq I'$ so that if $y\in \prod_{i'\in I'}R_{i'}'\Ball(M)$ satisfies $\|(a',b'c')-y\|_{2,\widetilde{F};X}<\eta$, then there is a weak$^{*}$-neighborhood $\mathcal{V}$ of $L_{y;X}$ with $\Gamma^{(n)}_{R'}(\mathcal{V})\subseteq \Gamma^{(n)}_{R'}(\mathcal{O})$.
Without loss of generality we may assume that $\eta<\varepsilon'$ and $\widetilde{F}\supseteq F'\cup G'$.  Since $C^{*}_{X}(a',b',c')\subseteq C^{*}_{X}(a,b,c)$. Apply Proposition \ref{prop: better KDT} to find a $Q\in \bC^{*}\ip{T_{i}:i\in I}^{I'}$ so that $\|Q(a,b,c)-(a',b',c')\|_{2,\widetilde{F}}<\eta$ and $\|Q_{i}\|_{R,\infty}\leq R_{i}'$ for all $i\in I'$.
Since $C^{*}_{X}(b)\supseteq C^{*}_{X}(b')$ we may, and will, assume that $Q|_{I_{2}'}\in \bC^{*}\ip{T_{i}:i\in I_{2}}^{I_{2}'}$. Choose a finite $G\subseteq I_{2}$ so that $Q_{i'}\in \bC^{*}\ip{T_{i}:i\in G}$ for all $i'\in G'$. Let $L=\|Q|_{G'}\|_{R,\Lip}$. Set  $\theta=\frac{\theta'}{L+1}$. Since 
\[\|Q|_{I_{1}'\sqcup I_{2}'}(a,b,c)-Q_{0}|_{I_{1}'\sqcup I_{2}'}(a,b,c)\|_{2,F';X}<2\varepsilon',\]
we have that 
\[\cU_{2}=\left\{\ell\in \Sigma_{R}:\sum_{i'\in F'}\|Q_{i'}(T')-Q_{0,i'}(T)\|_{L^{2}(\ell)}^{2}<4\varepsilon'^{2}\right\}\]
is a weak$^{*}$-neighborhood of $L_{a,b,c;X}$. 

% Note that $\|P(Q|_{I_{1}'\sqcup I_{2}'}(a,b))-(a,b)\|_{2,F;X}<2\varepsilon$. Thus
% \[\mathcal{U}_{0}=\left\{\ell\in \Sigma_{R}:\sum_{i\in F}\|P_{i}(Q|_{I_{1}'\sqcup I_{2}'}(T'))-T_{i}\|_{L^{2}(\ell)}^{2}<4\varepsilon^{2}\right\}\]
% is a neighborhood of $L_{a,b,c;X}$. 
By Lemma \ref{lem: key lem for invariance} (\ref{item: pushforward neighborhods}) and  (\ref{item: small perturbation neighborhood 2}), we may choose a neighborhood $\mathcal{U}_{3}$ of $L_{a,b,c;X}$ so that $Q(\Gamma^{(n)}(\mathcal{U}_{3}))\subseteq \Gamma^{(n)}(\mathcal{O})$ for all $n$. Let $\mathcal{U}=\mathcal{U}_{1}\cap \mathcal{U}_{2}\cap \mathcal{U}_{3}$. 
Suppose that $f\colon \Gamma^{(n)}(\mathcal{O})\to \Delta$ is $(\varepsilon',F'|\theta',G')$-injective. 
Suppose that $x,y\in \Gamma^{(n)}(\mathcal{U})$ and $f(Q(x))=f(Q(y))$ and that $\|x-y\|_{2,G}<\theta$. Then $\|Q(x)-Q(y)\|_{2,G'}\leq L\|x-y\|_{2,G}<\theta'$. Hence by choice of $f$ we have that $\|Q(x)-Q(y)\|_{2,F'}<\varepsilon'$. Our choice of $\mathcal{U}_{1},\mathcal{U}_{2}$ thus force:
\begin{align*}
\|x-y\|_{2,F}&\leq \varepsilon+\|P(Q_{0}|_{I_{1}'\sqcup I_{2'}}(x))-P(Q_{0}|_{I_{1}'\sqcup I_{2}'}(y))\|_{2,F}\\
&\leq 4\varepsilon+L\left[\|Q_{0}|_{I_{1}'\sqcup I_{2}'}(x)-Q|_{I_{1}'\sqcup I_{2}'}(x)\|_{2,F'}+\|Q_{0}|_{I_{1}'\sqcup I_{2}'}(y)-Q|_{I_{1}'\sqcup I_{2}'}(y)\|_{2,F'} \right]\\
&+\|P(Q|_{I_{1}'\sqcup I_{2}'}(x))-P(Q|_{I_{1}'\sqcup I_{2}'}(y))\|_{2,F}\\
&\leq 4\varepsilon+5L\varepsilon'<9\varepsilon.    
\end{align*}
Thus $f\circ Q$ is $(9\varepsilon,F|G,\theta)$-injective. So we have shown that 
\[\md_{9\varepsilon,F|\theta,G}(\mathcal{U}:c)\leq \md_{\varepsilon',F'|\theta',G'}(\mathcal{O}:c').\]
A fortiori,
\[\md_{9\varepsilon,F}(a|b:c)\leq \md_{\varepsilon',F'|\theta',G'}(\mathcal{O}:c').\]
Observe that $\varepsilon',F'$ depend upon $\varepsilon,F$ but do not depend upon $\mathcal{O},\theta',G'$. Moreover, $\theta',G'$ also do not depend upon $\mathcal{O}$. Thus we may fix $\varepsilon',F',\varepsilon,F,\theta',G'$ and take the infimum over $\mathcal{O}$ to obtain
\[\md_{9\varepsilon,F}(a|b:c)\leq \md_{\varepsilon',F'|\theta',G'}(a'|b':c').\]
Since $G',\theta'$ do not depend upon $\varepsilon,F,\varepsilon',F'$ and $\varepsilon,F$ do not depend upon $\varepsilon',F',\theta',G'$ we can take the infimum over $G',\theta'$ to see  that
\[\md_{9\varepsilon,F}(a|b:c)\leq \md_{\varepsilon',F'}(a'|b':c').\]
A fortiori,
\[\md_{9\varepsilon,F}(a|b:c)\leq \delta_{MD}(a'|b':c').\]
Taking the supremum over $\varepsilon,F$ shows that 
\[\delta_{MD}(a|b:c)\leq \delta_{MD}(a'|b':c').\]






\end{proof}

\end{document}
